\section{Related Works}

This section reviews the literature in relevant areas.

%\subsection{Foundation Models for Economic Forecasting}
%Recent advances in AI have seen the rise of \emph{foundation models}, predictive models that capture the dynamics of complex systems. Such world models aim to simulate or forecast the state of the environment, effectively allowing an AI to ``understand'' or predict future outcomes \cite{Ding2024}.  In economics, similar ideas have emerged under the framework of large pre-trained forecasting models.  \cite{Carriero2025} discuss the latest generation of time-series large language models (TSLMs) or time-series foundation models (TSFMs), which are large transformer-based models pretrained on diverse financial and economic data.  Examples include Google's TimesFM \cite{Das2023TimesFM}, Salesforce's Moirai (Woo et\ al.,~2024), and IBM's Tiny Time Mixers, all demonstrating that such models can outperform traditional forecasting baselines.  Notably, \cite{Zhu2025FinCast} introduce \emph{FinCast}, the first foundation model specifically designed for financial time-series, showing robust zero-shot forecasting performance and substantially improving on previous state-of-the-art methods.  These foundation models leverage vast pretraining data and can often be prompted or fine-tuned with a few examples (in-context learning) to adapt to new economic forecasting tasks \cite{Faw2025}.

%\subsection{Foundation models for economic forecasting}

%Recent work adapts the ``foundation model'' paradigm to forecasting by pretraining large neural networks on massive, heterogeneous time-series datasets and then transferring them across domains and tasks. In contrast to classical forecasting pipelines that fit task-specific models to limited panels, pretrained time-series foundation models aim to learn reusable representations of temporal dynamics that can be applied across datasets with minimal re-estimation. Representative approaches include tokenization-based probabilistic pretraining (e.g., Chronos) \citep{AnsariEtAl2024Chronos}, lag-augmented decoder-only transformers for probabilistic forecasting (Lag-Llama) \citep{RasulEtAl2023LagLlama}, and very large decoder-only models trained on web-scale time-series corpora (TimesFM) \citep{DasEtAl2024TimesFM}. Parallel to these ``pure time-series'' foundation models, an emerging line explores using large language models to extract signals from narrative and qualitative text and feed them into nowcasting systems, highlighting that relatively small but information-dense text sources can improve real-time forecasts \citep{deBondtSun2025ChatGPTNowcast,KwonEtAl2024BISLLMprimer}. At the same time, benchmarking work suggests that gains can be heterogeneous across variables and regimes, and that robustness to distribution shift and transparent evaluation against strong econometric baselines remain central for policy-facing deployment \citep{CarrieroEtAl2024MacroLLM}.

\subsection{Credit Risk Modeling in Financial Markets}
In traditional finance, credit risk modeling spans both individual-default models and systemic network approaches.  Classical structural models and reduced-form models quantify default probabilities of single entities, while network contagion models examine how defaults propagate through financial institutions.  For example, Dolfin et al. apply network epidemic models to credit contagion, highlighting how interconnected portfolios can amplify systemic credit risk \cite{math7080713}.  
In decentralized finance, credit risk manifests differently, as loans are typically over-collateralized and mediated by smart contracts. Bertomeu et al. propose simple aggregate risk measures for DeFi lending protocols, using only total deposits and borrowings, and report that systemic fragility surged around mid-2021 during crypto market turmoil \cite{BERTOMEU2024105321}. Conceptual studies have begun to map TradFi and DeFi risks together: Aufiero et al. review systemic risk mechanisms, finding that while basic risk types (leverage, liquidity shocks, correlated exposures) are common to both systems, DeFi's algorithmic execution and composability lead to unique propagation channels \cite{aufiero2025mappingmicroscopicsystemicrisks}.

\subsection{Machine Learning for Systemic Risk, Contagion, and Network Analysis}
Machine learning techniques are increasingly applied to systemic risk detection and contagion modeling \citep{BALMASEDA2023200240,gonon2025computingsystemicriskmeasures}. Graph neural networks (GNNs) in particular leverage network structure to improve risk prediction \citep{KipfWelling2017,VelickovicEtAl2018}. Balmaseda et al. show that a GNN-based model dramatically outperforms traditional ML for classifying systemic importance of banks in simulated networks \cite{BALMASEDA2023200240}. Gonon et al. integrate GNNs with explicit interbank liability networks, computing systemic risk measures as the minimum capital needed to secure the system \cite{gonon2025computingsystemicriskmeasures}. Their framework effectively learns contagion channels (e.g., Eisenberg-Noe clearing) by incorporating graph-structured data into the risk aggregation function. %These modern ML approaches build on network science foundations.  %For instance, \cite{math7080713} emphasize that network-theoretic methods are crucial for modeling correlated systemic credit risk.  
%Overall, machine-learning-driven network analysis enables end-to-end forecasting of systemic events, aiming to detect contagion pathways and provide early warnings for financial crises.

\subsection{Foundation Models in Finance and Economics}

Foundation models are increasingly penetrating into financial and economic domains by shifting the workflow from task-specific tools toward reusable representations learned from broad data. Instead of fitting a separate model for each target, recent approaches pretrain large neural networks, mostly using an architecture of transformer \citep{VaswaniEtAl2017Attention} and its variants, on massive data; and then transfer the pre-trained model across domains and tasks via prompting or light adaptation \citep{AnsariEtAl2024Chronos,RasulEtAl2023LagLlama,DasEtAl2024TimesFM}. In parallel, an emerging policy-facing literature uses large language models to extract quantitative signals from narrative text (e.g., reports, releases, and news) and incorporate them into nowcasting pipelines, suggesting that relatively small but information-dense text sources can improve real-time prediction when combined with standard indicators \citep{deBondtSun2025ChatGPTNowcast,KwonEtAl2024BISLLMprimer,CarrieroEtAl2024MacroLLM}. %Beyond ``pure'' time series, foundation-model architectures for structured data broaden the toolkit for financial applications: \fh{\cite{Ying2021Graphormer,HollmannEtAl2025TabPFN}}
These models are also adapted to graph inputs, which augment self-attention with graph structural encodings, demonstrating the feasibility of large pretrained models on networks that resemble webs of financial contracts or cross-asset relationships; a good example is Graphormer \citep{Ying2021Graphormer}. For tabular data, the Tabular Prior-data Fitted Network (TabPFN) \citep{HollmannEtAl2025TabPFN} is pretrained on millions of synthetic tabular tasks and reports strong out-of-the-box performance on diverse benchmarks. %with up to 10,000 samples, while also supporting generative uses such as density estimation and data generation. %Taken together, these advances imply that future DeFi risk models can integrate heterogeneous modalities; for example, coupling graph encoders for protocol exposure networks with transformer encoders over protocol- and token-level covariates, toward unified ``world models'' of financial ecosystems that forecast dynamics while aiming for robustness under regime shifts and measurement noise.


%\subsection{Graph and Tabular Foundation Models in Finance/DeFi}
%Emerging foundation-model architectures for graph and tabular data hold promise for financial applications.  Transformer-based models have been adapted to graph inputs: Graphormer \cite{Ying2021Graphormer} incorporates graph structural encodings (e.g., spatial and centrality features) into the Transformer, achieving state-of-the-art performance on a variety of graph prediction tasks.  This demonstrates the feasibility of large pretrained models on graph data, which could include networks of financial contracts or asset relationships.  For tabular data, recent work introduces powerful foundation models as well.  \cite{Hollmann2025} present the Tabular Prior-Data Fitted Network (TabPFN), a transformer model pretrained on millions of synthetic tabular tasks; TabPFN ``outperforms all previous methods'' on diverse tabular benchmarks with up to 10k samples.  As a generative model, TabPFN also supports fine-tuning, density estimation, and data generation, highlighting its versatility.  The development of such graph and tabular foundation models implies that future DeFi risk models can integrate heterogeneous data: for example, combining graph neural networks on protocol networks with transformer encoders on protocol-level features.  These advances provide a rich toolkit for building unified world models of financial ecosystems that operate across structured graph data and tabular features.
